\section{Construction of the prepared cycle}\label{sec:cycle}

We now combine the isolated boundary with the original log canonical
section.  The boundary supplies a divisor with torsion normal line;
the section supplies a canonical relation near its support.  A finite
cover removes the prime-to-$p$ part of the canonical exponent.  This
is the preparation of \cite[Sections~2.7--2.8]{Companion}, with the
support of the original boundary retained throughout the construction.

\begin{proposition}\label{prop:prepared}
Let $(X,B,N_0)$ be the data of Proposition~\ref{prop:reduction}
in the maximal-nef-dimension case.  There exist
a smooth projective threefold $V$, a very ample divisor $H_V$, and a
nonzero effective Cartier divisor
\[
 D=\sum_\alpha m_\alpha D_\alpha
\]
having all the properties in Theorem~\ref{thm:obstruction}.  In
particular, its support is connected with simple normal crossings,
$\gcd_\alpha(m_\alpha)=1$, and $\mathcal O_D(D)$ is torsion.
\end{proposition}

\begin{proof}
Equation~\eqref{eq:restriction-zero} and
Proposition~\ref{prop:separation} give the hypotheses of
Proposition~\ref{prop:isolation}. Choose its isolated model
$(Y,G_Y,N_Y,S_Y)$ and the reduced dlt modification
$f:(Y_d,\Sigma)\to(Y,S_Y)$ of Corollary~\ref{lem:dlt-boundary}.
Let $M$ be the positive Cartier multiple of $f^*S_Y$ supplied by
Proposition~\ref{prop:boundary-torsion}.

\emph{The canonical section and its support.}
Let $V_0$ be a smooth projective common resolution of $X$ and $Y_d$,
chosen to dominate the intervening models.  Write
\[
 q:V_0\longrightarrow X,\qquad h_0:V_0\longrightarrow Y_d,
 \qquad r=f\circ h_0:V_0\longrightarrow Y.
\]
These models share the open set
\[
 U_X=X_{\mathrm{reg}}\setminus
       (\operatorname{Supp}N_0\cup P_{\mathrm{out}}),
\]
where $P_{\mathrm{out}}$ is the union of the boundary components
disjoint from $\operatorname{Supp}N_0$.  We choose the common
resolution to be an isomorphism over this open.  Projective resolution
and resolution of divisor supports in dimension three permit this
choice \cite{CP,CJS}.  The pullback equalities through the partial
MMP give
\begin{equation}\label{eq:cycle-common-pullbacks}
 q^*N_0=r^*N_Y.
\end{equation}

The integer $m$ was chosen so that
$N_0\sim m(K_X+B)$ and both $mK_X$ and $mB$ are Cartier.  It gives a
rational section $\theta$ of $\omega_{V_0}^{\otimes m}$ with
\begin{equation}\label{eq:log-canonical-section}
 \operatorname{div}(\theta)
   =q^*N_0+m(K_{V_0}-q^*K_X)-m q^*B.
\end{equation}
Indeed, the right side is an integral divisor linearly equivalent to
$mK_{V_0}$.  Put $E_0=\operatorname{Supp}q^*N_0$.  Every prime
component in \eqref{eq:log-canonical-section} is either contained in
$E_0$ or disjoint from it.  For $q^*B$ this follows from boundary
separation on $X$.  For a $q$-exceptional divisor, its irreducible
center lies in
\[
 X\setminus U_X
   =\operatorname{Supp}N_0\cup P_{\mathrm{out}}
       \cup\operatorname{Sing}X.
\]
A terminal threefold is regular in codimension two, so
$\operatorname{Sing}X$ is finite
\cite[Lemma~2.3]{Companion}.  The center is therefore contained in
$\operatorname{Supp}N_0$, or is disjoint from it.  Taking inverse
images gives the assertion for the exceptional divisor.

Let $F_0$ be the union of the components of
$\operatorname{div}(\theta)$ outside $E_0$.  Thus
\begin{equation}\label{eq:cover-closed-set}
 Z_0:=E_0\cup\operatorname{Supp}\operatorname{div}(\theta)
          =E_0\sqcup F_0.
\end{equation}
The disjointness is the point of this description.  Components of
$q^*N_0$ can cancel in \eqref{eq:log-canonical-section}; we retain all
of $E_0$ when choosing the open set preserved by the subsequent
resolution.

\emph{Removing the prime-to-$p$ exponent.}
Write $m=p^b m'$ with $b\geq0$ and $(m',p)=1$.  Choose a nonzero
rational generator $\eta$ of $\omega_{V_0}^{\otimes p^b}$ and write
$\theta=u\eta^{m'}$, where $u\in k(V_0)^\times$.  Adjoin one root
$v$ of $v^{m'}=u$ and normalize $V_0$ in the resulting field.
This is a finite separable cover, possibly of degree smaller than
$m'$.  Over $V_0\setminus\operatorname{Supp}\operatorname{div}(\theta)$,
it is an integral component of the torsor of $m'$th roots of the
invertible section $\theta$.  This description is independent of the
rational generator $\eta$.  Since $m'$ is invertible in $k$, that
torsor is finite \'etale.

Resolve the normalization and the total inverse image of $Z_0$,
preserving the regular \'etale open over $V_0\setminus Z_0$.  We
obtain a smooth projective threefold and a generically finite
separable morphism
$\pi:V\to V_0$ for which the total pullback support of $N_0$ is
simple normal crossings.  The morphisms used to compare divisors are
\[
\begin{tikzcd}[column sep=large,row sep=large]
 V\arrow[r,"\pi"] &V_0\arrow[r,"q"]\arrow[d,"h_0"']&X\\
 &Y_d\arrow[r,"f"']&Y.
\end{tikzcd}
\]
Here $q,h_0,f$ are birational; $\pi$ is generically finite and
separable.  Set $\rho=q\circ\pi$ and $h=h_0\circ\pi$.  By
\eqref{eq:cycle-common-pullbacks}, $\rho^*N_0=(f\circ h)^*N_Y$,
while the Cartier divisor supplied by
Proposition~\ref{prop:boundary-torsion} pulls back as $h^*M$.

The rational section $v\eta$ is a section of
$\pi^*\omega_{V_0}^{\otimes p^b}$.  The determinant of the
differential is a nonzero section of the rational line
$\omega_V\otimes\pi^*\omega_{V_0}^{-1}$; denote it by $J_\pi$.
Its nonvanishing at the generic point follows from separability.
Consequently
\begin{equation}\label{eq:canonical-root-section}
 \psi=(v\eta)J_\pi^{p^b}
       \quad\text{is a rational section of }
                \omega_V^{\otimes p^b}.
\end{equation}
Both factors are invertible on $\pi^{-1}(V_0\setminus Z_0)$: the first is
a root of an invertible section, and the second is the Jacobian of
an \'etale morphism.  Thus the divisor of $\psi$ is supported on
$\pi^{-1}(Z_0)$.  By \eqref{eq:cover-closed-set}, each of its prime
components lies in $\operatorname{Supp}\rho^*N_0$ or is disjoint
from that support.

\emph{A connected primitive Cartier divisor.}
Put $M_V=h^*M$ and choose a connected component $C$ of
$\operatorname{Supp}M_V$.  Since $M$ is a positive multiple of
$f^*S_Y$, and $S_Y$ is disjoint from the other components of
$N_Y$, the set $C$ is also a connected component of
$\operatorname{Supp}\rho^*N_0$.  Write the part of $M_V$ supported
on $C$ as
\begin{equation}\label{eq:primitive-cycle}
 aD,\qquad D=\sum_\alpha m_\alpha D_\alpha,
 \qquad a=\gcd_{D_\alpha\subset C}
                       \operatorname{mult}_{D_\alpha}(M_V).
\end{equation}
The divisor $D$ is Cartier because $V$ is smooth, and its
multiplicities are positive with gcd one.  Its support is connected
and simple normal crossings.

Every curve in $\operatorname{Supp}M_V$ maps under $h$ to a point or
to a curve in $\Sigma$.  The projection formula and
Proposition~\ref{prop:boundary-torsion} therefore show that $M_V$
is numerically trivial on every component of its support.  Near $C$
we have $M_V=aD$, so $D|_{D_\alpha}\equiv0$ for every $\alpha$.
For a curve not contained in $\operatorname{Supp}D$, effectivity
gives nonnegative intersection with $D$.  For a curve in its support,
the intersection is zero.  Thus $D$ is nef, and every very ample
divisor $H_V$ satisfies
\begin{equation}\label{eq:primitive-intersections}
 D H_V^2>0,\qquad D^2H_V=0.
\end{equation}

Torsion also survives on the entire Cartier scheme.  Choose $t>0$
with $\mathcal O_M(tM)\simeq\mathcal O_M$.  The pullback divisor
$M_V$ is scheme-theoretically $V\times_{Y_d}M$: its local equation
is the pullback of a local equation of $M$, which is a nonzero
divisor since $V$ is integral and $h$ is dominant.  Pulling back
the trivialization gives
$\mathcal O_{M_V}(tM_V)\simeq\mathcal O_{M_V}$; no flatness
assumption on $h$ is needed.  Restricting along $D\subset M_V$,
and using $M_V=aD$ near $C$, gives
\begin{equation}\label{eq:primitive-torsion}
 \mathcal O_D(taD)\simeq\mathcal O_D.
\end{equation}
Hence $\mathcal O_D(D)$ is torsion.

The section $\psi$ of \eqref{eq:canonical-root-section} has, on
a neighborhood $W$ of $C$, divisor supported on $D$.  Indeed, its
other components are disjoint from $C$ and can be removed from
that neighborhood.  Thus there are integers $b_\alpha$ such that
\begin{equation}\label{eq:prepared-canonical-relation}
 \omega_V^{\otimes p^b}|_W
     \simeq\mathcal O_W\Bigl(\sum_\alpha b_\alpha D_\alpha\Bigr).
\end{equation}
\emph{Numerical proportionality and very general points.}
It remains to show that no $D$-trivial curve passes through a very
general point of $V$.  Put $N_V=\rho^*N_0$.  It is a nonzero nef
divisor, numerically trivial on every component of its support by
the projection formula.  Since $\operatorname{Supp}D$ is contained
in that support, we have
\[
 N_VH_V^2>0,\qquad N_V^2H_V=N_VDH_V=0.
\]
Together with \eqref{eq:primitive-intersections}, the Hodge index
theorem gives
\begin{equation}\label{eq:cycle-proportionality}
 D\equiv cN_V,\qquad
 c=\frac{DH_V^2}{N_VH_V^2}>0.
\end{equation}
We verify that this is numerical equivalence on every curve, rather
than only on a fixed hyperplane section, following
\cite[Section~2.8]{Companion}.

Let $\Gamma\subset V$ be any integral curve.  For $l\gg0$, there
is an integral surface $T\in|lH_V|$ containing $\Gamma$.  To see
this, twist the ideal of $\Gamma$ sufficiently far and then multiply
its generators by very ample sections.  The resulting linear system
has no fixed divisorial component and gives a birational map onto
its image off $\Gamma$.  Bertini irreducibility gives an irreducible
general member, which is generically reduced on that birational open.
As a Cartier divisor on smooth $V$, it has no embedded components,
and is therefore integral.
On a smooth projective resolution $\widetilde T\to T$,
let $d,n,e$ be the pullbacks of $D|_T,N_V|_T,H_V|_T$, respectively.
Their intersections are
\[
 d^2=n^2=dn=0,\qquad e^2>0,\qquad
 de=lDH_V^2,\quad ne=lN_VH_V^2.
\]
Hence $d-cn$ has square zero and is orthogonal to $e$.  The
Hodge index form is negative definite on the orthogonal complement
of the positive-square class $e$, so $d-cn$ is numerically zero.
The projective surjection $\widetilde T\to T$ supplies an integral
curve dominating $\Gamma$: take the closure of a closed point of
the fiber over its generic point.  The projection formula then gives
$(D-cN_V)\cdot\Gamma=0$.  This proves
\eqref{eq:cycle-proportionality}.

Finally let $Z_i\subset X$ be the countably many proper closed
subsets excluded by the nef reduction of $N_0$.  Remove from $V$
the non-quasi-finite locus of $\rho$ and all the sets $\rho^{-1}(Z_i)$.
These are proper closed subsets, since $\rho$ is dominant and
generically finite.  Their complement contains very general points
because $k$ is uncountable.  Every curve through any point of this
complement maps to a curve on $X$.  If it were $D$-trivial, then
\eqref{eq:cycle-proportionality} and the projection formula would
make its image $N_0$-trivial, contradicting maximal nef dimension.
All properties in Theorem~\ref{thm:obstruction} have now been
established.
\end{proof}

\begin{proof}[Proof of Theorem~\ref{thm:main}]
Theorem~\ref{thm:obstruction} contradicts
Proposition~\ref{prop:prepared}.  The maximal-nef-dimension case is
therefore impossible.  By Proposition~\ref{prop:reduction}, it remains
to apply abundance in nef dimension at most two
\cite[Theorem~1.8]{XuNonvanishing}.  This gives semiampleness, and
faithful descent along the algebraically closed field extension gives
it over the original field.
\end{proof}
